\subsection{A five-coordinate counteralgorithm}

Let $\widehat y:=D^{-1/2}\widehat x$ and define
\begin{equation}
 z_i:=\frac{2(6-i)(7-i)}{11}
 \quad(1\leq i\leq5),
 \qquad
 \widehat y_i:=
 \begin{cases}
  \alpha_n z_i,&1\leq i\leq5,\\
  0,&i\geq6.
 \end{cases}
 \label{eq:five-site-certificate-candidate}
\end{equation}
Thus
$(z_1,z_2,z_3,z_4,z_5)=(60,40,24,12,4)/11$ and
$\widehat x=D^{1/2}\widehat y$ has exactly five nonzero coordinates.

\begin{theorem}[A degree-four supported polynomial beats the product scale;
\ProvedStatus]
\label{thm:five-site-supported-prefix-counteralgorithm}
For every $n\geq8$, the vector in
\cref{eq:five-site-certificate-candidate} satisfies the strict version of
\cref{eq:sparse-path-residual-task}.  It belongs to
$\mathcal K_5(Q_n,b)$ and hence is returned by an admissible degree-four
$\mathsf{CertPrefixPoly}$ execution with resource vector
\begin{equation}
 \mathcal C_{\rm CPP}=\bigl(
 \Theta(1),\Theta(1),1,0,\Theta(1),\Theta(1),0,
 \Theta(1),\Theta(1),\Theta(1),5\bigr).
 \label{eq:certificate-prefix-poly-resource-vector}
\end{equation}
On the identical task and adjacency interface, a fully charged five-row
append-only $LDL^\top$ response emits the same vector with
\begin{equation}
 \mathcal C_{\rm cert\text{-}LDL}=\bigl(
 \Theta(1),\Theta(1),1,0,\Theta(1),0,\Theta(1),
 \Theta(1),\Theta(1),0,5\bigr).
 \label{eq:certificate-path-response-resource-vector}
\end{equation}
Both algorithms are deterministic and use only exact algebraic cells.
Their realized final exposed set may be taken as $[6]$, of degree volume
$\nu_{\rm fin}=11$.  Consequently neither
$\Omega(\nu_{\rm fin}/\sqrt{\alpha_n})=\Omega(n)$ nor
$\Omega(\nu_n/\sqrt{\alpha_n})=\Omega(n^2)$ lower-bounds the charged work of
this named sparse-certificate recurrence class on the endpoint path.
\end{theorem}

\begin{proof}
Put $a=(1+\alpha_n)/2$ and $\eta=(1-\alpha_n)/2$.  Because the first six
vertices are not the far endpoint when $n\geq8$, direct substitution in the
ambient-degree path rows gives
\begin{equation}
 \frac{1}{\alpha_n}
 \bigl[D^{-1/2}(b-Q_n\widehat x)\bigr]_i
 =
 \begin{cases}
  (1-50\alpha_n)/11,&i=1,\\
  (1-41\alpha_n)/11,&i=2,\\
  (1-25\alpha_n)/11,&i=3,\\
  (1-13\alpha_n)/11,&i=4,\\
  (1-5\alpha_n)/11,&i=5,\\
  (1-\alpha_n)/11,&i=6,\\
  0,&i\geq7.
 \end{cases}
 \label{eq:five-site-exact-residual}
\end{equation}
For example, at $\alpha_n=0$ the scaled row equations reduce to
\[
 1-\frac{z_1-z_2}{2}
 =\frac{z_{i-1}-2z_i+z_{i+1}}4
 =\frac{z_5}{4}=\frac1{11},
\]
where the middle identity uses $2\leq i\leq5$ and $z_6=0$.
The $\alpha_n$ terms give the coefficients displayed in
\cref{eq:five-site-exact-residual}.  Since
$\alpha_n\leq1/64$, all six displayed values are positive and at most
$1/11<1/10$.  This proves the strict certificate.

The endpoint Krylov triangularity established in the companion exact-CG
theorem gives
\[
 \mathcal K_5(Q_n,b)
 =\operatorname{span}\{e_1,\ldots,e_5\}.
\]
Hence there is a polynomial $p$ of degree at most four with
$\widehat x=p(Q_n)b$.  A concrete supported realization forms
$w_0=b$ and $w_{j+1}=Q_nw_j$ for $0\leq j<4$, then solves the resulting
constant $5\times5$ lower-triangular Krylov coordinate system and emits the
linear combination equal to $\widehat x$.  The four row actions scan prefixes
$[1],[2],[3],[4]$, of total ambient degree volume
$1+3+5+7=16$.  The coefficient solve and vector combination are constant
exact recurrence arithmetic, not an uncharged response.  Exposing and caching
the first six path rows, the four supported products, intermediate cells, and
the exact final support-plus-boundary verifier all use constant work and
space.  There is no preprocessing or response object, and exactly five final
cells are emitted.  This proves every coordinate in
\cref{eq:certificate-prefix-poly-resource-vector}; its constant adjacency
round count is an implementation charge, not an adaptivity lower bound.

For the response arm, append the first five tridiagonal pivots, multipliers,
and forward messages.  Equivalently, solve the leading five-row system with
right-hand side equal to the first five coordinates of
$b-D^{1/2}g$, where
$g:=D^{-1/2}(b-Q_n\widehat x)$ is given by
\cref{eq:five-site-exact-residual}.  One reverse pass gives exactly
$\widehat x_{1:5}$; row six supplies the boundary residual.  Charge every
adjacency reply and verifier operation to $C_{\rm adj}$ or $C_{\rm ctl}$,
every pivot, message, and reverse application to $C_{\rm resp}$, every stored
or scratch cell to its memory coordinate, and the five output cells to
$C_{\rm emit}$.  No intermediate full vector is maintained, so
$C_{\rm mat}=0$.  All quantities are constant, proving
\cref{eq:certificate-path-response-resource-vector} on the same task and
interface.

Finally, $\nu_n=2(n-1)$, $\nu_{\rm fin}=\vol([6])=11$, and
$1/\sqrt{\alpha_n}=n$, whereas both charged executions have constant
nonmemory work.
\end{proof}

\paragraph{Exact scope of the refutation.}
The theorem refutes a product lower bound based only on supported-prefix
polynomial generation plus the actual sparse residual certificate.  It does
not refute the full-vector exact
$\mathsf{DiagSpecPoly}(r)$ obstruction in
\cref{thm:surviving-path-polynomial-obstruction}: the two results have
different output tasks.  It also does not shorten the ordinary-CG trajectory
proved in the companion note.  Ordinary CG annihilates the old-prefix
residual and leaves a frontier value just above the threshold; the
counteralgorithm instead spends the permitted residual slack across six local
rows.  The result is path-specific, tied to $\alpha=n^{-2}$ and
$\epsppr=1/10$, and has no finite-precision or graph-uniform consequence.
